% (c) Nikita Lisitsa, lisyarus@gmail.com, 2025

\documentclass[10pt]{beamer}

\usepackage[T2A]{fontenc}
\usepackage[russian]{babel}
\usepackage{minted}

\usepackage[most]{tcolorbox}
\tcbuselibrary{minted}

\usepackage{graphicx}
\graphicspath{ {./images/} }

\usetheme{metropolis}

\usepackage{adjustbox}

\usepackage{color}
\usepackage{soul}

\usepackage{hyperref}

\definecolor{codebg}{RGB}{29,35,49}

\setminted{fontsize=\footnotesize}

\setlength\partopsep{-\topsep}

\definecolor{blue}{rgb}{0,0,1}
\definecolor{red}{rgb}{1,0,0}

\usemintedstyle{tango}

\makeatletter
\newcommand{\slideimage}[1]{
  \begin{figure}
    \begin{adjustbox}{width=\textwidth, totalheight=\textheight-2\baselineskip-2\baselineskip,keepaspectratio}
      \includegraphics{#1}
    \end{adjustbox}
  \end{figure}
}
\makeatother

\title{Компьютерная графика}
\subtitle{Домашнее задание 1: Планетная система}
\date{2026}

\setbeamertemplate{footline}[frame number]

\begin{document}

\frame{\titlepage}

\begin{frame}[fragile]
\slideimage{screenshot.png}
\end{frame}

\begin{frame}[fragile]
\frametitle{Задание}
\begin{itemize}
\item Сцена с Солнцем, планетами, и звёздами
\item Вокруг некоторых планет есть спутники или кольца
\item Планеты вращаются вокруг Солнца, спутники -- вокруг планет
\item Звёзды неподвижны и ``бесконечно далеки''
\item Планеты и Солнце крутятся вокруг своей оси и как-нибудь раскрашены
\item Цвета у Солнца меняются анимированно, вокруг Солнца -- фейковое сияние
\end{itemize}
\end{frame}

\begin{frame}[fragile,label=sun_planets]
\frametitle{Камера}
\begin{itemize}
\item Камера должна достаточно свободно двигаться
\item Идеально: WASD -- двигаемся вперёд/назад/влево/вправо, движение мыши -- вращаем камеру
\item Удобно хранить 3D координаты самой камеры + два угла поворота (надо разобраться, вокруг каких осей)
\item В начале программы сделайте \mintinline{cpp}|app.grab_mouse(true)| -- это не даст курсору мыши выйти за пределы окна
\end{itemize}
\end{frame}

\begin{frame}[fragile,label=sun_planets]
\frametitle{Солнце и планеты}
\begin{itemize}
\item Нужно сгенерировать кодом 3D модель сферы, вариантов много -- UV-sphere, icosphere, subdivided cube
\item Главное -- чтобы она была с индексами, без дырок, и без дублирующихся вершин (\textit{manifold mesh})
\item На все планеты и Солнце -- один буфер с 3D моделью сферы, но у каждого свой буфер с цветами вершин (как-нибудь их сгенерируйте)
\item Шейдер Солнца просто выводит цвета как есть, а шейдер планет использует освещение по формуле
\usemintedstyle{lightbulb}
\begin{tcblisting}{listing engine=minted, minted language=wgsl, minted options={fontsize=\scriptsize}, listing only, colback=codebg, colframe=codebg, rounded corners}
// Здесь position - позиция в 3D пространстве
// Считаем, что Солнце находится в начале координат
// ambient - что-нибудь в духе 0.01..0.1
color * (ambient + max(0.0, dot(normal, -normalize(position))))
\end{tcblisting}
\usemintedstyle{tango}
\item Матрицу \mintinline{cpp}|model| нужно применять к вектору нормали
\item Цвета Солнца как-нибудь анимируются на CPU и каждый кадр загружаются на GPU
\end{itemize}
\end{frame}

\begin{frame}[fragile]
\frametitle{3D модель сферы}
\slideimage{sphere-combined.png}
\end{frame}

\begin{frame}[fragile,label=sun_planets]
\frametitle{Солнце и планеты}
\begin{itemize}
\item Планеты движутся вокруг Солнца и вращаются вокруг своей оси (ось может быть наклонена)
\item Спутники движутся вокруг своих планет и вращаются вокруг своей оси
\item Солнце тоже вращается вокруг своей оси
\item Все эти повороты, сдвиги, и масштабирования должны делаться одной матрицей \mintinline{cpp}|model|
\end{itemize}
\end{frame}

\begin{frame}[fragile]
\slideimage{moons.png}
\end{frame}

\begin{frame}[fragile,label=sun_planets]
\frametitle{Кольца}
\begin{itemize}
\item Для колец нужно сгенерировать отдельную 3D модель ``кольца''
\item Она будет односторонней (без внутренности), и будет неправильно освещаться с ``обратной'' стороны
\item Для этого есть полезный флаг во фрагментном шейдере \mintinline{wgsl}|@builtin(front_facing) : bool| -- если он равен \mintinline{rust}|false|, домножим вектор нормали на -1
\end{itemize}
\end{frame}

\begin{frame}[fragile]
\slideimage{rings.png}
\end{frame}

\begin{frame}[fragile,label=sun_planets]
\frametitle{Свечение солнца}
\begin{itemize}
\item Для свечения Солнца можно использовать ту же 3D модель кольца, но заставить её повернуться в сторону камеры
\item Например, взять направление из Солнца на камеру за ось Y, взять любые две перпендикулярных ей и друг другу оси X и Z, и построить из них матрицу поворота
\item Для динамического свечения можно тоже динамически обновлять буфер с цветами
\end{itemize}
\end{frame}

\begin{frame}[fragile]
\slideimage{sun.png}
\end{frame}

\begin{frame}[fragile,label=sun_planets]
\frametitle{Звёзды}
\begin{itemize}
\item Для звёзд лучше использовать модель попроще (например, октаэдр или низкополигональную сферу)
\item Все звёзды должны рисоваться одной командой и использовать instancing: нужен отдельный буфер с инстансами звёзд (например, положение, цвет, размер), и правильная настройка render pipeline
\item Цвета звёзд можно взять случайными, а можно взять из формул blackbody spectrum :)
\item Позиции звёзд можно сделать случайными единичными векторами
\item Чтобы звёзды казались бесконечно далеко, рисуем их самыми первыми в кадре, без теста глубины и без записи в буфер глубины (это может быть отдельный render pass, или правильно настроенный pipeline), и используем для них особую view-матрицу камеры, в которой есть поворот, но нет сдвига на позицию камеры
\end{itemize}
\end{frame}

\begin{frame}[fragile]
\slideimage{stars.png}
\end{frame}

\begin{frame}[fragile,label=sun_planets]
\frametitle{Баллы}
\footnotesize
\begin{itemize}
\item \textbf{1 балл}: рисуются хоть какие-то сферы
\item \textbf{1 балл}: модель сферы -- manifold mesh
\item \textbf{1 балл}: камера двигается по WASD
\item \textbf{1 балл}: камера крутится мышью
\item \textbf{1 балл}: есть центральное Солнце и планеты вокруг него
\item \textbf{1 балл}: есть спутники
\item \textbf{1 балл}: есть кольца
\item \textbf{1 балл}: все вращаются вокруг своей оси
\item \textbf{2 балла}: у Солнца анимируются цвета
\item \textbf{1 балл}: есть звёзды
\item \textbf{2 балла}: все звёзды рисуются одним instanced draw
\item \textbf{2 балла}: У Солнца есть ``свечение''
\end{itemize}
Всего: \textbf{15 баллов}

Защита заданий на практике \textbf{через 2 недели}.
\end{frame}

\end{document}
